\section{\texorpdfstring{$J$-全纯曲线的整体性质}{J-全纯曲线的整体性质}}
\label{sec:global-properties-j-holomorphic-curves}

格罗莫夫在 (Gr85) 中引入的一个富有创见、并彻底改变了辛拓扑领域的想法，是研究殆复映射空间
\[
 f:(\Sigma,j)\longrightarrow(M,J),
\]
其中 $(\Sigma,j)$ 是另一个殆复流形。

\begin{definition}\label{def:j-holomorphic-map}
殆复映射 $f:(\Sigma,j)\to(M,J)$ 称为一个 $(j,J)$-全纯映射。
\end{definition}

不过，除非 $j$ 和 $J$ 都可积，否则当 $\dim_{\mathbb R}\Sigma>2$ 时，甚至这类映射的局部存在定理也不可能成立。另一方面，当 $\dim_{\mathbb R}\Sigma=2$ 时，有下面的局部存在定理。我们经常略去 $j$，或者把一对 $(j,f)$ 称为一个\emph{$J$-全纯映射}，并把它的像称为一条（无参数化的）$J$-全纯曲线。\footnote{译者注：原书在定理题头与 (NiW63) 书目中均将合著者 William B. Woolf 的姓误作 Wolf；据原论文（DOI 10.2307/1970126），该结果由阿尔伯特·奈恩黑斯（Albert Nijenhuis）与威廉·伍尔夫（William B. Woolf）证明。}

\begin{theorem}[奈恩黑斯--伍尔夫，(NiW63)]\label{thm:nijenhuis-wolf-local-disc}
对每一点 $p\in M$ 以及 $x$ 处的每个复切 $2$-平面 $\tau_p$，都存在 $\epsilon>0$ 和一个 $J$-全纯圆盘
\[
 f:(D^2(\epsilon),j)\longrightarrow(M,J),
\]
使得 $f(0)=p$ 且 $\tau_p=\operatorname{Image}T_0f$。
\end{theorem}

这个定理的证明依赖于非线性偏微分方程的椭圆理论；注意，殆复条件
$df\circ j=J\circ df$ 给出一个拟线性一阶椭圆偏微分方程。事实上，对 $(D^2(\epsilon),j)$ 的情形，它等价于
\[
 \frac{\partial f}{\partial x}
 +J(f)\frac{\partial f}{\partial y}=0,
\]
其中
\[
 \frac{\partial f}{\partial x}:=df\left(\frac\partial{\partial x}\right),
 \qquad
 \frac{\partial f}{\partial y}:=df\left(\frac\partial{\partial y}\right),
\]
而 $(x,y)$ 是 $D^2(\epsilon)$ 的标准坐标。当然，当 $J$ 可积时，如果把 $J$ 与 $M$ 的复坐标 $(z_1,\ldots,z_n)$ 上乘以 $\sqrt{-1}$ 等同起来，并把 $\Sigma$ 上的 $j$ 与对 $z=x+\sqrt{-1}y$ 乘以 $\sqrt{-1}$ 等同起来，这个方程就恰好成为全纯曲线方程。

现在研究 $J$-全纯曲线的一些整体性质。为了与后面的计算记号保持一致，我们将从记号 $Tu$ 改用 $du$，并把它看成一个取值于 $u^*TM$ 的向量值一形式。于是有分解
\begin{equation}
 du=\partial_Ju+\overline\partial_Ju,
 \label{eq:du-j-complex-decomposition}
\end{equation}
其中 $\partial_Ju$（相应地，$\overline\partial_Ju$）是 $du$ 关于 $(j,J)$ 的复线性（相应地，复反线性）部分：
\begin{align}
 \partial_Ju&:=\frac12(du-J\circ du\circ j),
 \label{eq:j-complex-linear-part}\\
 \overline\partial_Ju&:=\frac12(du+J\circ du\circ j).
 \label{eq:j-antilinear-part}
\end{align}
于是，$u$ 是 $J$-全纯的，当且仅当 $\overline\partial_Ju=0$。

在 $(\Sigma,j)$ 上固定一个埃尔米特度量 $h$。映射
\[
 du:(T\Sigma,h)\longrightarrow(TM,g)
\]
关于任意黎曼度量 $g$ 的范数 $|du|$，定义为
\[
 |du|_g^2:=\sum_{i=1}^2|du(e_i)|_g^2,
\]
其中 $\{e_1,e_2\}$ 是 $T\Sigma$ 关于 $h$ 的一个正交归一标架。这个定义与正交归一标架 $\{e_1,e_2\}$ 的选择无关。用坐标表示时，设 $(x^1,x^2)$ 是 $\Sigma$ 上的坐标，$(y_1,y_2,\ldots,y_{2n})$ 是 $M$ 上的坐标，并写
\[
 g=\sum_{\alpha,\beta}g_{\alpha\beta}\,dy^\alpha dy^\beta,
 \qquad
 h=\sum_{i,j}h_{ij}\,dx^i dx^j,
 \qquad (h^{ij})=(h_{ij})^{-1}.
\]
于是称为\emph{调和能量密度函数}的 $|du|_g^2$ 由下式给出：
\[
 |du|_g^2=\sum_{i,j,\alpha,\beta}
 g_{\alpha\beta}(u(x))h^{ij}(x)
 \frac{\partial u^\alpha}{\partial x^i}
 \frac{\partial u^\beta}{\partial x^j}.
\]

下述结论来自 $J$-全纯映射的定义以及 $J$ 与 $\omega$ 的相容性。

\begin{proposition}\label{prop:j-holomorphic-energy-identities}
记 $g_J=\omega(\cdot,J\cdot)$，并把相伴范数记为
$|\cdot|=|\cdot|_J$。固定 $(\Sigma,j)$ 的一个埃尔米特度量 $h$，并考虑光滑映射 $u:\Sigma\to M$。则有
\begin{enumerate}[label=(\arabic*)]
 \item $|du|^2=|\partial_Ju|^2+|\overline\partial_Ju|^2$；
 \item $2u^*\omega=(-|\overline\partial_Ju|^2+|\partial_Ju|^2)dA$，其中 $dA$ 是 $\Sigma$ 上度量 $h$ 的面积形式。
\end{enumerate}
\end{proposition}

\begin{proof}
上面给出的范数 $|du|^2$ 之定义与正交归一标架 $\{e_1,e_2\}$ 的选择无关。因此，在 $T\Sigma$ 上选取关于 $h$ 的正交归一局部标架，使得 $je_1=e_2$，并以 $\{\alpha_1,\alpha_2\}$ 表示其对偶标架。

依定义，
\[
 |du|^2=|du(e_1)|^2+|du(e_2)|^2
\]
以及
\[
 |\overline\partial_Ju|^2
 =|\overline\partial_Ju(e_1)|^2+|\overline\partial_Ju(e_2)|^2,
 \qquad
 |\partial_Ju|^2
 =|\partial_Ju(e_1)|^2+|\partial_Ju(e_2)|^2.
\]
由 $e_2=je_1$ 的选择，
\[
 \overline\partial_Ju(e_1)
 =\frac{du(e_1)+Jdu(e_2)}2,
 \qquad
 \overline\partial_Ju(e_2)
 =\frac{du(e_2)-Jdu(e_1)}2,
\]
并且
\[
 \partial_Ju(e_1)
 =\frac{du(e_1)-Jdu(e_2)}2,
 \qquad
 \partial_Ju(e_2)
 =\frac{du(e_2)+Jdu(e_1)}2.
\]
利用自同态 $J$ 与 $j$ 分别关于度量 $g$ 与 $h$ 正交这一事实，作直接计算便立即得到恒等式 (1)。

为了证明恒等式 (2)，把面积形式写成
\[
 dA=\alpha_1\wedge\alpha_2.
\]
由 $J$ 与 $\omega$ 的相容性，直接计算给出
\[
\begin{split}
 |\partial_Ju|^2-|\overline\partial_Ju|^2
 &=-2\langle du(e_1),Jdu(e_2)\rangle
 =2\langle Jdu(e_1),du(e_2)\rangle\\
 &=2\omega(du(e_1),du(e_2))
 =2u^*\omega(e_1,e_2).
\end{split}
\]
这就蕴含 (2)，证明完毕。
\end{proof}

从这些恒等式得到 $J$-全纯曲线关于辛形式 $\omega$ 的下述\emph{校准}性质。
\stfhIndex{I071}{jiao^zhun^xing^zhi}{校准性质（calibrated）}

\begin{corollary}\label{cor:calibrated-energy-inequality}
对任意光滑映射 $u:\Sigma\to M$，都有
\[
 \frac12\int_\Sigma|du|_J^2\,dA\geq\int_\Sigma u^*\omega,
\]
并且等号成立恰好当 $u$ 是 $J$-全纯的。
\end{corollary}

注意，当 $\Sigma$ 是无边界闭曲面时，$\int u^*\omega$ 在一个固定同调类中是常数。因此，任何 $J$-全纯曲线都使固定同调类中的调和能量
$\frac12\int_\Sigma|du|_J^2dA$ 达到最小。这立即给出下面的结论。

\begin{proposition}\label{prop:j-holomorphic-area-homology}
当 $u:(\Sigma,j)\to(M,J)$ 是 $J$-全纯映射时，有
\[
 \operatorname{Area}_{g_J}u
 =\frac12\int|du|_J^2=\int u^*\omega.
\]
因此，$\operatorname{Area}_{g_J}u$ 只依赖于 $u$ 所代表的同调类。特别地，对代表固定同调类 $A=[u]$ 的任意 $J$-全纯映射
$u:(\Sigma,j)\to(P,J)$，
\begin{equation}
 \operatorname{Area}_{g_J}(u)=[\omega](A).
 \label{eq:j-holomorphic-area-homology-class}
\end{equation}
\end{proposition}

\begin{corollary}\label{cor:j-holomorphic-image-minimal-surface}
在 $u$ 于 $\Sigma$ 上的每个正则点附近，$J$-全纯映射的像关于度量
$g_J=\omega(\cdot,J\cdot)$ 是极小曲面。
\end{corollary}

\begin{remark}\label{rem:varying-j-uniform-area-estimate}
在应用中，必须改变殆复结构，并考虑包含于 $\mathcal J_\omega$ 的某个紧子集 $K$ 中的一族 $J$。不过，在估计映射序列 $\{f_i\}$ 并建立其收敛论证时，我们必须在相容度量中使用给定的固定度量 $g$。于是，\eqref{eq:j-holomorphic-area-homology-class}应替换为
\[
 \operatorname{Area}_g(f)\leq C(A,K),
\]
其中 $f$ 是任意 $J$-全纯映射，$J\in K\subset\mathcal J_\omega$，而 $C(A,K)$ 是只依赖于 $A$ 与 $K$ 的常数。这个一致面积估计，将是获得一致 $C^1$ 估计、进而研究 $J$-全纯映射之紧性性质（其中 $J$ 在一个紧集中变化）的重要一步。
\end{remark}
